\justifying
\NSFCBodyText

\subsubsubsection{面向物流状态的可靠性自适应多模态表征}

现有方法多将同步完整输入作为默认条件，或对缺失模态统一填补。本项目把时间偏移、传感质量和跨模态一致性共同纳入时变可靠性估计，使振动、倾斜、电子封志、光感和包装图像对事件判断的贡献随观测条件动态变化，同时约束运输事件识别、包装损伤定位和概率校准。其创新性由模态删除、时间错位和噪声扰动下的事件级性能、mAP/IoU、Brier/ECE 及退化曲线共同检验。

\subsubsubsection{风险信息价值驱动的通信--计算--能量协同机制}

现有弱网研究常分别优化上传、卸载、压缩或能耗。本项目在同一端侧状态中描述事件风险、网络、缓存、算力、电量与断点确认信息，把本地推理、摘要/压缩、离线缓存、断点续传、优先上传和动态工作周期作为联合动作，研究高风险证据按时送达与资源消耗之间的条件边界。该机制不以提出新通信协议为目标，而以相同事件流和网络轨迹下的高风险漏报、P95 时延、传输字节和单位任务能耗进行反证。

\subsubsubsection{不同时标仓储证据的风险演化与可校准分级预警}

区别于温湿度、烟雾、视频和门禁各自设阈值，本项目研究缓慢环境变化与突发安防事件的跨源共现、持续性和可靠性变化，形成可追溯的风险状态、预警等级及响应建议。通过单源/独立阈值对照、来源删除、时间置乱和完整仓储事件外层测试，验证复合风险是否获得更早且更可信的识别；运输阶段状态仅作为可选先验，不构成仓储主线的成立前提。
